\chapter{Finding a Quasioptimal N-Cyclic Arbitrage Walk with MCTS based REINFORCE}
In chapter we go through the RL framework \& neural network model used
to find profitable walks in the CFMM arbitrage problem and present results.
Firstly we introduce the model for the underlying MDP, states, actions and
rewards. The policy is parametrized  via the deep GNNs of the line graph of
the problem and the training setup uses the setup from chapter \ref{sec:
gen}.

\section{Agent and Environment}
The framework the agent is learning to find a walk in a graph that does not
use repeated edges (trail), starts and ends at the same node, i.e. is
circular, and is profitable. The graph is represented by ERC20 \cite{erc20}
tokens (referred to as assets) as nodes and decentralized exchange pools,
CFMM with two assets, as edges (referred to as pools). Each pool has a swap
fee, which usually is in the range of $0.3\%-1\%$, and a marginal exchange
rate. The marginal exchange rate, unique to the pool, is inversely
bidirectional. Meaning, while swapping token $0$ for token $1$ in the pool
the usual marginal exchange rate is used, but when swapping from token $1$ to
token $0$ the inverse marginal exchange rate is used. In both cases the
referred term is marginal exchange rate. `Used' here is meant in the context
of the learning setup, in the actual setting the automated market maker
formula is used to determine the amount exchanged \ref{chap: cfmms}.
Throughout we will use the term `pool' and edge, `asset' and node
interchangeably.
\newline

The agent starts from a given node (asset) and chooses actions from a set of
valid actions. Each action represents a trade, where the agent swaps an asset
for another asset through a specific pool. The valid actions rule out the
edges that have already been used. The idea behind this is the following: If
there is an arbitrage opportunity, and the trades are conducted in a
chronological order, at each swap the pool's `capacity' will be exerted and
the pools' exchange rate will be in balance with the overall market. Hence, the
pool has no liquidity for an arbitrage to be used again.
\newline

When the agent arrives at a terminal state, there are two possibilities. The
first possibility is when the agent has arrived back to the starting point,
then the product of the exchange rates is calculated to give the profit of
the walk. In this case, at each trade a trading penalty can utilized to
account for the fees of each transaction. The second possibility is when the
agent arrives at a dead end, meaning the agent has not reached the starting
point and has no valid actions left.
\newline

To bring it all together the MDP is based on a multi bidirectional graph $G =
(V, E, W)$, where $V$ is the set of nodes (assets), $E$ the set of edges
(pools) and $W$ the set of edge weights (prices based on the specific pool).
An edge $e_{ij}^{k}  \in E$ represents the pool $k$ trading assets $i$ $j$,
it has the weight
\begin{align}
w_{ij}^{k} =\frac{p_i^{k} }{p_j^{k} },
\end{align}
where $p_i^{k} $ and $p_j^{k} $ are the marginal exchange rates of the assets
$i$ and $j$ in pool $k$ respectively as defined in equation \ref{eq:
marginal_rate}. The edge $e_{ji}^{k}$ has the weight
\begin{align}
w_{ji}^{k}  = \frac{1}{w_{ij}^{k} } = \frac{p_j^{k}}{p_i^{k}}.
\end{align}
The index $k$ denotes the $k$-th pool trading asset $i$ and asset $j$, since
there could be multiple pools trading the same asset. The underlying MDP is a
$4$ tuple $(S, A, P, R)$, deterministic, and is specified by
:\newline
\textbf{Actions}. The action space $A$ are all the possible actions the agent
can take on the graph defined above. The agent acts by taking an action $a_t
\in A$ at step $t$ by selecting an edge $e_{ij}^{k}$ leading to adjacent nodes
and thereby transitioning from node $i$ to node $j$ at pool $k$. The
action is valid if the edge $e_{ij}^{k}$ has not been selected throughout the
episode.\newline
\textbf{States}. The state space $S$, and a state $s_t \in S$ is the edge set
traversed by the agent up to step $t$ of the episode. The state $s_t$ is
terminal if there are no valid actions or the agent has reached the origin
node. If the agent has not reached a terminal state then the agent selects an
action $a_t \in A$ and transitions to the next state $s_{t+1} = s_t \cup a_t$
recursively.\newline
\textbf{Transition}. The transition probabilities $P$, where
$\mathbb{P}[s'|s, a] = 1$ is the probability of landing in state $s'$ by
taking action $a$ at state $s$. The MDP is deterministic.\newline
\textbf{Rewards}. The reward probabilities $R$, since the MDP is
deterministic the reward function at state $s_t$ is $r(s_t)$. Throughout the
MDP, the agent is rewarded $r(s_t) =0$, except when landing back at the
origin node. All together the agent is rewarded by
\begin{align}
    r(s_t) =
    \begin{cases}
        \max\left\{0, \sum_{w_{ij}^{k}  \in s_t}\ln(w_{ij}^{k} \cdot \tau)\right\}  \quad
        &\text{terminal state and reached the
        origin node}\\
            0 \quad &\text{else},
    \end{cases}
\end{align}
where $w_{ij}^{k}$ are the weights of the walk of $s_T$, i.e. marginal
exchange rates, $\tau \in[0, 1]$ is the trade penalty which incorporates a
rough estimate of the transaction cost for each trade. Throughout the
training we set $\tau = 1$, however in production it might be useful to set
an adaptive value to get an estimate if the overall trade will be profitable
after paying the protocol-transaction fees. If the sum of the logarithm of
the marginal exchange rates is bigger then zero then the trade is profitable
and otherwise it is not. The natural logarithm comes here as a big
convenience to characterize profitable and nonprofitable trades,it also
avoids $256$-bit integer overflows, (most of the blockchains do not use
floating point number, but unsigned $256$-bit integers) which needs special
attention when coding in production. In summary the agent is only rewarded if
end node matches the start node.

\section{Policy Parametrization}
We are going to parametrize the policy $\pi_\theta$, such that given a state
$s$ the output of $\pi_\theta(s)$ will be a probability vector of size $n = |E|$.
This probability vector will be masked based on the valid actions at state
$s$ and then normalized. To archive this we will encode the state $s$ through
the line graph $L(G)$ of the graph used in the MDP. The nodes of the line
graph $L(G)$ represent the edges of the graph $G$. Two nodes are connected by
an edge in $L(G)$ if and only if their corresponding edges in $G$ share a
common endpoint in $G$. Instead of considering both edges $e_{ij}^{k}$,
$e_{ji}^{k}$, we will only consider the corresponding pool connection as an
edge. A node of the line graph $L(G)$ will have the following attributes
\begin{itemize}[nosep]
    \item Natural logarithm of the marginal exchange rate corresponding to
        the pool
    \item Binary entry stating if the pool was used $1$ or not used $0$
    \item Binary entry stating if agent is currently using asset $0$ of the
        pool, $1$ or not
        $0$.
    \item Binary entry stating if agent is currently using asset $1$ of the
        pool, $1$ or not
        $0$.
\end{itemize}
Given a state $s$ the input to the policy parametrization will be encoded
through graph as stated in the list above. The architecture of the deep GNN
parametrization of the policy is as follows. Given the node attributes, the
encoder computes the node embeddings $h_i^{0}$ for all nodes $i$ through a
learned liner projection to dimension $d_h = 256$. The embeddings are then
updated via. $N$ Residual Multihead Attention Layers \cite{gatv2}. Each layer
consists of two sublayers, a multihead attention layer (GAT) with $16$ heads,
for message passing between the nodes which is flattened by a linear
connection back to the original dimension $d_h$. And a node-wise fully
connected feed-forward layer (FF) with $1024$ hidden dimension. Both
sublayers have a Batch Norm (BN) and a residual connection similar to
\cite{rl_coopt}. The following is a formalization of the residual block
\begin{align}
    \hat{h}_{i} &= \text{BN}^{l}\left( h_i^{(l-1)}
    + \text{GAT}_i^{l}(h_1^{(l-1)},\ldots,h_n^{(l-1)})\right)\\
    h_i^{(l)} &= \text{BN}^{l}\left( \hat{h}_{i} + \text{FF}_i^{l}(\hat{h}_i) \right).
\end{align}
After which we will do a contextual embedding by concatenating the graph
embedding $\hat{h}^{(N)} = \frac{1}{n}\sum_{i=1}^{n} h_i^{(N)}$, the node
embedding of the initial state $h^{\pi(s_0)}_i$ and the vectors $h_i^{N}$ to
the vector
\begin{align}
    h^{(c)}_i = \hat{h}^{(N)} \|  h^{\pi(s_0)}_i \| h_i^{(N)},
\end{align}
where $\|$ is the concatenation operator. The contextual embeddings
$h^{(c)}_i$ are passed into the policy head. The policy head consists of
GAT layer with $16$ heads passed to a Batch Norm layer and ultimately
compressed down to an output dimension of $2$ per node through a linear
connection with Relu activation function
\begin{align}
    \pi_\theta(s) = \text{softmax}\left(   W_p \cdot
    \text{relu}\left(\text{BN}_p(\text{GAT}_p(h^{(c)}_1,\ldots,h^{(c)}_n))\right)\right),
\end{align}
where $W_p \in \mathbb{R}^{2\times d_h}$ is the learned linear connection,
and the index $p$ denotes the distinct weights used for the policy head. To
get the policy vector we flatten and pass the softmax at the end.


\section{Algorithm and Training}
The training is purely unsupervised based on the MCTS and REINFORCE algorithm
as described in Section \ref{sec: gen}. At each self play The marginal
exchange rates of the prices are chosen randomly, some block number $b$ on
the Ethereum blockchain. The next state is chosen based on $N_{\text{MCTS}}$
searches with selection probabilities computed through $\pi_\theta$. In each
iteration the agent plays the network $N_{\text{play}}$ times until reaching
a terminal state. Upon reaching a terminal state: the state, the action, the
reward and the policy vector of the MCTS algorithm are saved. After the self
play stage the model parameters are updated based on the self play data. We
alter the Rollout-Simulation pahse of the MCTS algorithm by instead of doing
one random rollout we do $N_{\text{sim}}$ number of rollouts. We fully
utilize the REINFORCE algorithm updates through \ref{alg: reinforce}, where
the loss function can be written as
\begin{align}
    L(\theta) = -\frac{1}{N_{\text{batch}}}
    \sum_{i=1}^{N_{\text{batch}}}G_{i} y_i
    \ln \pi_{\theta} (s^{b_i}_{i_t}),
\end{align}
where $G_i$ is the sum of the rewards as in the REINFORCE algorithm, and
$y_i$ a one hot vector with $1$ at the index of the action taken at the end
of an MCTS. The action taken at the end of an MCTS is determined by the MCTS
improved policy $\pi^{\text{MCTS}}$. The $Q$ values are estimated through the
MCTS and parametrized policy $\pi_\theta$ search. The full reinforcement
learning algorithm algorithm in pseudo-code is the following
\begin{algorithm}[H]
\caption{RL Algorithm}\label{alg: rl_alg}
    \begin{algorithmic}
        \While{within computational budget}
        \For {$i=1,\ldots, N_{\text{play}}$}
            \State Initialize state $s_0 = s$ at block $b$
            \While{not terminal}
            \State $\pi^{MCTS}(s)  \gets \text{MCTS}(s)$
            \State  $s \gets \argmax_{a}\pi^{MCTS}(a|s) $
            \State save $s$, $\pi^{\text{MCTS}}(s)$ and $a$
            \EndWhile
        \EndFor
        \State $\theta \gets Adam(\nabla L(\theta))$
        \EndWhile
    \end{algorithmic}
\end{algorithm}
The code can be found at \cite[Github]{github}.
\newline

The weighs of the model are updated using the Adam optimizer \cite{adam} with
a learning rate of $\alpha = 10^{-4} $, a weight decay of $\lambda =
10^{-4}$, forgetting factors $\beta_1 = 0.9$, $\beta_2 = 0.999$ and
$\varepsilon = 10^{-8}$ to prevent division by zero. We set the discount rate
constant $\gamma = 0.98$ from REINFORCE \ref{alg: reinforce}. We use
$N_{\text{batch} }=25$, $N_{\text{play}} = 200$, $N_{\text{MCTS}} = 100$
with $N_{\text{sim}}=100$ (MCTS simulations) with $100$ iterations. The
marginal exchange rates are pulled from the block range $t \in
[18\cdot10^{6}, 21.8\cdot10^6]$. The number of residual block layers is
$N=20$, these do not share weights. The UCB exploration-exploitation
parameter $C$ in the MCTS algorithm for choosing the next child in the tree
is set to $C=1.4$.
\section{Data}
The data used is publicly available on the Ethereum mainnet blockchain and
can be accessed through an archive node query. The pools used are all of
UniswapV2 \cite{uniswap} and UniswapV3 \cite{uniswapv3_paper}. The difference
between V2 pools and V3 pools is that the V3 pools have concentrated
liquidity pools, but both the follow constant product function. At the
current time of writing there are around $4.3\cdot10^{5}$ pools and
approximately $4.1\cdot10^{5}$ different ERC20 \cite{erc20} tokens in these
pools. To narrow down the meaningful pools, firstly we use pools that have
been created before block $18\cdot10^{6}$ and secondly we use a degree
filter. Assets which have a smaller degree then $4$ are filtered out. The
question still remains if these pools are still in use currently. Hence, the
pools with low liquidity are also filtered out, by looking at their two token
reserves, the $\text{ETH}$ value of the reserves is crucial here. If the
reserve is worth less then $100\; \text{ETH} $ it is filtered out. We also
look at the gradient of the marginal exchange rates for all the pools, and
filter out the pools with low price change. For the training process the
marginal exchange rates of the pools are extracted from block $18\cdot10^{6}$
until the current block every $3600$ blocks which is equal to around every
$12$ hours ($\sim12s$ is the time until the next block). During each
self-play, a random block is sampled from which the marginal exchange rates
are taken. Each state starts with the agent holding the WETH asset, from
which the agent chooses a pool to conduct a swap with WETH to another token
untill reaching a terminal state.
\section{Result}
Average reward on the dataset per iteration
IMAGE HERE
Calculations were performed using supercomputer resources provided by the Vienna Scientific Cluster (VSC)
